% ==============================================================================
% SEÇÃO 5.3: RESULTADOS PRELIMINARES E AUDITORIA DE EQUIDADE INTERSECCIONAL
% Proposta de Qualificação de Doutorado - PPGC/UFRGS
% ==============================================================================

\section{Resultados Preliminares e Auditoria de Equidade Interseccional}
\label{sec:resultados_preliminares}

Nesta seção, apresentam-se os resultados empíricos da avaliação de desempenho preditivo e de equidade algorítmica interseccional obtidos com os modelos de referência (\textit{baselines}). A análise estrutura-se a partir dos experimentos executados sobre três bases de dados com diferentes contextos sociodemográficos e epidemiológicos: \textbf{Adult} (US Census), \textbf{COMPAS} (Recidivismo Criminal) e \textbf{SINASC} (DATASUS, Saúde Neonatal no Brasil).

Os experimentos foram conduzidos sob o protocolo rigoroso de Validação Cruzada Aninhada (\textit{Nested Cross-Validation}), composto por 5 \textit{folds} externos para estimação de generalização e 3 \textit{folds} internos para calibração de hiperparâmetros, contemplando múltiplos objetivos de otimização ($\text{accuracy}$, $\text{recall}$, $\text{precision}$, $\text{roc\_auc}$ e $\text{pr\_auc}$) aplicados aos algoritmos \textit{Random Forest} (RF) e \textit{Gradient Boosting} (GB).

\subsection{Desempenho Preditivo Global e Análise Comparativa}
\label{subsec:desempenho_global}

A Tabela~\ref{tab:modelos_resumo} sintetiza o comportamento preditivo e os indicadores de disparidade interseccional agregados em nível de base de dados e arquitetura. As métricas reportam a média e o desvio-padrão computados sobre os \textit{folds} de teste externos ao longo de todas as funções objetivo avaliadas no \textit{inner-CV}.

\begin{table}[htbp]
\centering
\small
\caption{Desempenho preditivo e disparidade interseccional consolidada (médias $\pm$ desvios-padrão entre as métricas de otimização no \textit{inner-CV}).}
\label{tab:modelos_resumo}
\begin{tabular}{llccccc}
\toprule
\textbf{Dataset} & \textbf{Modelo} & \textbf{Acurácia} & \textbf{ROC-AUC} & \textbf{PR-AUC} & \textbf{Max AAOD} $\downarrow$ & \textbf{Sens. Gap} $\downarrow$ \\
\midrule
Adult & Gradient Boosting & 0.860 $\pm$ 0.008 & 0.915 $\pm$ 0.009 & 0.813 $\pm$ 0.018 & 0.0936 $\pm$ 0.0260 & 0.1247 $\pm$ 0.0547 \\
Adult & Random Forest     & 0.843 $\pm$ 0.014 & 0.899 $\pm$ 0.010 & 0.778 $\pm$ 0.022 & 0.1155 $\pm$ 0.0196 & 0.1658 $\pm$ 0.0544 \\
\midrule
COMPAS & Gradient Boosting & 0.638 $\pm$ 0.003 & 0.691 $\pm$ 0.000 & 0.698 $\pm$ 0.000 & 0.2614 $\pm$ 0.0076 & 0.2303 $\pm$ 0.0055 \\
COMPAS & Random Forest     & 0.641 $\pm$ 0.001 & 0.690 $\pm$ 0.000 & 0.697 $\pm$ 0.000 & 0.2644 $\pm$ 0.0075 & 0.2354 $\pm$ 0.0039 \\
\midrule
SINASC & Gradient Boosting & 0.906 $\pm$ 0.000 & 0.532 $\pm$ 0.002 & 0.913 $\pm$ 0.001 & 0.0000 $\pm$ 0.0000 & 0.0000 $\pm$ 0.0000 \\
SINASC & Random Forest     & 0.906 $\pm$ 0.000 & 0.532 $\pm$ 0.000 & 0.914 $\pm$ 0.000 & 0.0000 $\pm$ 0.0000 & 0.0000 $\pm$ 0.0000 \\
\bottomrule
\multicolumn{7}{l}{\footnotesize $\downarrow$ Indica métricas onde valores menores representam maior equidade interseccional.} \\
\end{tabular}
\end{table}

A distribuição global das métricas de desempenho preditivo é ilustrada na Figura~\ref{fig:overall_performance}, comparando o comportamento dos modelos em termos de Acurácia e Revocação (\textit{Recall}). Os \textit{boxplots} em escala de cinza capturam a variabilidade induzida pelas diferentes métricas de otimização interna, enquanto os marcadores mapeiam conjuntamente a base de dados (código de cores) e o critério de calibração de hiperparâmetros (formato do marcador).

\begin{figure}[htbp]
\centering
\includegraphics[width=0.95\linewidth]{fig_overall_performance.pdf}
\caption{Desempenho global dos modelos em termos de Acurácia e Revocação. Cores representam os \textit{datasets}, e as formas dos marcadores indicam a função objetivo otimizada no \textit{inner-CV}.}
\label{fig:overall_performance}
\end{figure}

No \textit{dataset} Adult, o algoritmo \textit{Gradient Boosting} demonstrou superioridade consistente frente ao \textit{Random Forest}, atingindo ROC-AUC médio de $0.915 \pm 0.009$ (com pico de $0.919$ quando otimizado explicitamente por $\text{roc\_auc}$) e PR-AUC de $0.813 \pm 0.018$, contra $0.899 \pm 0.010$ e $0.778 \pm 0.022$ do \textit{Random Forest}. Essa vantagem estende-se à métrica de acurácia média ($0.860$ vs $0.843$).

No \textit{dataset} COMPAS, ambos os modelos apresentaram desempenho estatisticamente equivalente, com ROC-AUC estabilizado em torno de $0.691$ e acurácia próxima a $64\%$. Observa-se baixíssima dispersão entre os \textit{folds} e entre as funções de custo (desvios-padrão $\le 0.003$), refletindo o limite empírico de separabilidade estatística das variáveis disponíveis no cadastro penal de Broward County.

\subsection{Impacto das Funções Objetivo de Otimização no Inner-CV}
\label{subsec:impacto_funcoes_objetivo}

A análise comparativa entre as funções objetivo utilizadas na seleção interna de hiperparâmetros revelou que a escolha da métrica de calibração afeta diretamente a severidade do viés algorítmico, mesmo em modelos sem intervenção explícita de \textit{fairness}.

No \textit{dataset} Adult:
\begin{enumerate}
    \item \textbf{Otimização por Acurácia e Recall:} Favoreceu o equilíbrio operacional do classificador, atingindo um $\text{Max AAOD}$ significativamente menor: $0.0736$ (otimizado por acurácia) e $0.0813$ (otimizado por \textit{recall}), com $\text{Sensitivity Gap}$ restrito a $0.0873$ e $0.0982$, respectivamente.
    \item \textbf{Otimização por Precisão:} Ao penalizar severamente os falsos positivos no \textit{inner-CV}, o algoritmo adotou limiares conservadores de classificação, concentrando as predições favoráveis ($>50\text{k}$) exclusivamente em instâncias de alta certeza. Esse conservadorismo amplificou drasticamente a injustiça interseccional: o $\text{Sensitivity Gap}$ explodiu para $0.2209$ no \textit{Gradient Boosting} e $0.2163$ no \textit{Random Forest}, com o $\text{Max AAOD}$ saltando para $0.1389$.
\end{enumerate}

No \textit{dataset} COMPAS, a variação da função objetivo provocou variações no balanço entre sensibilidade e especificidade, com a otimização por $\text{pr\_auc}$ alcançando a maior revocação global ($0.724$ no GB), porém preservando patamares alarmantes de disparidade de chances igualadas ($\text{Max AAOD}$ entre $0.2512$ e $0.2699$).

Esses achados corroboram a hipótese central da tese: o ajuste de modelos baseado estritamente em métricas cegas de desempenho é insuficiente e pode agravar desproporcionalmente a penalização de subgrupos historicamente vulneráveis.

\subsection{Auditoria Interseccional Pós-Treinamento e Análise de Subgrupos}
\label{subsec:auditoria_interseccional}

A Figura~\ref{fig:fairness_distribution} expõe a relação de compromisso (\textit{trade-off}) entre a capacidade discriminativa global ($\text{ROC-AUC}$) e a máxima injustiça interseccional observada ($\text{Max AAOD}$). 

\begin{figure}[htbp]
\centering
\includegraphics[width=0.95\linewidth]{fig_fairness_distribution.pdf}
\caption{Distribuição de ROC-AUC versus Máxima Injustiça Interseccional ($\text{Max AAOD}$).}
\label{fig:fairness_distribution}
\end{figure}

Formalmente, a métrica $\text{Max AAOD}$ mensura a pior violação de chances igualadas (\textit{Equalized Odds}) entre qualquer subgrupo interseccional $g \in \mathcal{G}$ e o grupo de referência $g_{\text{ref}}$:
\begin{equation}
\text{Max AAOD} = \max_{g \in \mathcal{G}} \frac{1}{2} \left( |\text{TPR}_g - \text{TPR}_{g_{\text{ref}}}| + |\text{FPR}_g - \text{FPR}_{g_{\text{ref}}}| \right)
\end{equation}
onde $\mathcal{G} = \mathcal{A} \times \mathcal{S}$ define a partição de subgrupos no espaço cartesiano dos atributos sensíveis (e.g., Gênero $\times$ Raça).

A desagregação dos resultados nos subgrupos interseccionais revela disparidades profundas que permaneceriam ocultas sob auditorias unidimensionais convencionais:

\subsubsection{Adult Dataset: A Penalização Composta de Mulheres Negras}
Na base Adult, o grupo de referência (\textit{Male \& White}, $N=8.884$) apresentou uma taxa de predição favorável de $24.8\%$, com $\text{TPR} = 59.3\%$ e $\text{FPR} = 7.9\%$. Em contrapartida:
\begin{itemize}
    \item O subgrupo \textbf{Female \& Black} ($N=717$) obteve taxa de predição favorável de apenas \textbf{$3.8\%$}, correspondendo a um Impacto Disparatado pós-treino ($\text{Post-DI}$) de \textbf{$0.149$} e $\text{AAOD} = 0.1020$, com a taxa de verdadeiros positivos caindo para $45.9\%$.
    \item Mulheres brancas (\textit{Female \& White}, $N=3.867$) registraram taxa favorável de $7.8\%$ ($\text{Post-DI} = 0.306$).
    \item Homens negros (\textit{Male \& Black}, $N=732$) registraram taxa favorável de $12.7\%$ ($\text{Post-DI} = 0.507$).
\end{itemize}

\noindent\textbf{Constatação Teórica:} Uma auditoria univariada de gênero concluiria genericamente que mulheres recebem menos predições favoráveis. Todavia, a auditoria interseccional revela que mulheres negras sofrem uma penalização composta: sua chance de classificação favorável é \textbf{6,5 vezes menor} que a de homens brancos e menos da metade da observada para mulheres brancas. Esse achado valida a tese de interseccionalidade de Kimberl\'e Crenshaw e os resultados empíricos de \cite{buolamwini2018gender}.

\subsubsection{COMPAS: Disparidade Racial e de Gênero no Risco Penal}
No COMPAS, onde a classe favorável $Y=1$ corresponde à ausência de reincidência (baixo risco), os dados evidenciam que:
\begin{itemize}
    \item \textbf{Male \& African-American} ($N=1.007$): registrou a menor taxa de classificação favorável entre todos os grupos ($44.8\%$, $\text{Post-DI} = 0.585$), a menor taxa de verdadeiros positivos ($\text{TPR} = 58.1\%$) e o maior desvio absoluto de chances igualadas ($\text{AAOD} = 0.2629$).
    \item Em contrapartida, \textbf{Female \& Caucasian} ($N=187$) obteve taxa favorável de $72.0\%$ ($\text{Post-DI} = 0.943$), $\text{TPR} = 81.4\%$ e $\text{AAOD} = 0.0436$.
    \item \textbf{Male \& Caucasian} ($N=611$) atingiu taxa favorável de $64.1\%$ e $\text{TPR} = 73.8\%$.
\end{itemize}

O $\text{Sensitivity Gap}$ superior a \textbf{$23.5$ pontos percentuais} entre homens afro-americanos e indivíduos caucasianos comprova que os modelos \textit{baseline} absorvem o viés do sistema judiciário, penalizando sistematicamente réus negros com taxas desproporcionais de falsos positivos desfavoráveis.

\subsection{Comportamento sob Desbalanceamento Severo: O Caso do SINASC/DATASUS}
\label{subsec:sinasc_desbalanceamento}

O \textit{dataset} SINASC ($N = 2.474.535$ registros de nascidos vivos no Brasil) representa o desafio de aplicação de equidade em saúde pública em larga escala. O objetivo clínico consiste em prever o baixo peso ao nascer ($<2.500\text{g}$, classe minoritária $Y=0$, representando $9.44\%$ dos casos), em oposição ao peso adequado ($\ge 2.500\text{g}$, classe majoritária $Y=1$, representando $90.56\%$). Os atributos sensíveis combinam Raça/Cor da Mãe (5 categorias) e Faixa Etária Materna (3 categorias), totalizando 15 subgrupos interseccionais.

Os resultados da Tabela~\ref{tab:modelos_resumo} expõem o fenômeno denominado \textbf{Colapso Trivial por Desbalanceamento}:
\begin{enumerate}
    \item Ambos os algoritmos convergiram para uma acurácia nominal de \textbf{$90.57\%$}, rigorosamente igual à prevalência da classe majoritária.
    \item A revocação da classe majoritária foi de $1.000$, porém a discriminabilidade real permaneceu próxima ao acaso ($\text{ROC-AUC} \approx 0.532 \pm 0.002$).
    \item As métricas de equidade convergiram para valores aparentemente nulos: $\text{Max AAOD} = 0.0000$, $\text{Sensitivity Gap} = 0.0000$ e $\text{Post-DI} = 1.000$ em todos os 15 subgrupos.
\end{enumerate}

\noindent\textbf{Diagnóstico Científico:} Essa equidade aparente é uma \textbf{ilusão estatística induzida pela degeneração do modelo}. Ao predizer invariavelmente a classe majoritária para $100\%$ das parturientes, as taxas de acerto e de erro falso tornam-se constantes ($\text{TPR}_g = 1.0$ e $\text{FPR}_g = 1.0$) para todas as combinações de raça e idade:
\begin{equation}
|\text{TPR}_g - \text{TPR}_{g_{\text{ref}}}| = |1.0 - 1.0| = 0 \quad \text{e} \quad |\text{FPR}_g - \text{FPR}_{g_{\text{ref}}}| = |1.0 - 1.0| = 0 \implies \text{AAOD}_g = 0
\end{equation}

Esse comportamento empírico sustenta três contribuições para a tese:
\begin{itemize}
    \item Comprova que modelos convencionais falham catastroficamente em bases epidemiológicas sem estratégias de mitigação sensíveis a desbalanceamento.
    \item Demonstra que métricas de paridade podem ser enganosas quando desprovidas de capacidade preditiva real ($\text{PR-AUC}$).
    \item Justifica formalmente a adoção de um algoritmo evolucionário multiobjetivo para guiar o aprendizado para regiões não-degeneradas do espaço de soluções.
\end{itemize}

\subsection{Síntese Crítica e Justificativa para a Arquitetura Metodológica Proposta}
\label{subsec:sintese_justificativa}

Os experimentos preliminares fornecem o suporte empírico necessário para a arquitetura metodológica proposta nesta qualificação:

\begin{enumerate}
    \item \textbf{Rede Neural Multitarefa com Camada Reversa de Gradiente (GRL):} \\
    Para desacoplar a informação dos múltiplos atributos protegidos sem impor suposições restritivas de independência estatística, propõe-se uma arquitetura onde um codificador latente compartilhado $E(X)$ minimiza a perda da tarefa principal $L_{\text{task}}$ enquanto ramos adversariais buscam predizer os atributos sensíveis $S_1$ e $S_2$:
    \begin{equation}
    L_{\text{total}} = L_{\text{task}} - \left( \lambda_1 L_{\text{adv1}} + \lambda_2 L_{\text{adv2}} \right)
    \end{equation}
    A inversão do gradiente via GRL força a extração de representações latentes invariantes a subgrupos interseccionais, neutralizando a raiz do viés observado no Adult e no COMPAS.

    \item \textbf{Otimização Multiobjetivo Evolucionária (NSGA-II):} \\
    Os experimentos evidenciaram que o espaço de calibração entre desempenho e equidade contém descontinuidades e platôs triviais (conforme visto no SINASC). Em vez de fixar arbitrariamente pesos $\lambda_1, \lambda_2$, propõe-se utilizar o algoritmo \textbf{NSGA-II} \cite{deb2002fast} para mapear a \textbf{Fronteira de Pareto empírica}:
    \begin{equation}
    \max_{\lambda_1, \lambda_2} \; \left[ f_1 = \text{PR-AUC}, \quad f_2 = -\text{Max AAOD}, \quad f_3 = -\text{Sensitivity Gap} \right]
    \end{equation}
    permitindo identificar modelos operacionais viáveis sem colapso trivial.

    \item \textbf{Auditoria e Interpretabilidade Pós-Processamento com Valores de Shapley (SHAP):} \\
    A integração de explicações por valores de Shapley desagregadas por subgrupos \cite{lundberg2017unified} permitirá auditar se o mecanismo adversarial eliminou o uso de atributos substitutos (\textit{proxy features}), garantindo conformidade e auditabilidade no contexto de políticas públicas de saúde.
\end{enumerate}

Em conclusão, os dados preliminares estabelecem não apenas a viabilidade técnica da pesquisa, mas também a sua necessidade teórica e social, alinhando a proposta aos mais elevados padrões científicos de aprendizado de máquina confiável e justo.
